\section{Contraste de las dos pruebas DRG.3 con la demostración de Hash\_DRBG}
\label{sec:hash-comparison}
El §1073 de AIS v3 dicta una conclusión para Hash\_DRBG y remite a un paso
distinto para cada requisito. El §998 reconoce el ejemplo AES/SHA; el §1007
encarga la prueba de CTR al solicitante. Para comparar sus argumentos hay que
hacer coincidir el problema de predicción, el momento de compromiso y el
horizonte. En los dos perfiles de este trabajo hay un número interno de
128 bits por petición. Un lector que agregue varios bloques ejecuta varias
transiciones completas. La vista, la transición y el presupuesto de seguridad
de esos perfiles se especifican en las secciones~\ref{sec:ais-lean},
\ref{sec:ctr-verified} y~\ref{sec:ctr-finite-reductions}
\cite[§§998,1007,1033--1075]{ais}.

\paragraph{Qué pasos se trasladan del ejemplo del BSI.}
En §§1035--1038, BSI escribe toda la salida de una petición de Hash\_DRBG
como evaluaciones en puntos consecutivos, con posible truncamiento (5.56),
y calcula el siguiente estado (5.58). Los §§1039--1046 motivan hipótesis
para Hash y sus transformaciones específicas; (5.65), por ejemplo, trata
evaluaciones en puntos con diferencias conocidas. En §§1048--1052 se concede
información adicional al adversario: diferencias de estados, bloques completos
y una ventana más larga. Los §§1053--1064 estudian si el argumento hash del bloque
buscado coincide con uno previo. La cota \(2^{-330}\) de (5.78) limita ese
\emph{evento} bajo el horizonte y las hipótesis del ejemplo. Para las dos tuplas
presentes se debe calcular la probabilidad de predecir bajo sus propias
transiciones. La decisión de usar una petición de un bloque elimina los
casos de truncamiento y de objetivos interiores a una petición.

\begin{center}\small
\begin{tabularx}{\linewidth}{@{}lYYY@{}}
\toprule
DRG.3 & Paso en Hash\_DRBG & AES/SHA & CTR directo\\
\midrule
.1, .4 & §§1031--1033 y 1074--1075. El origen y la entropía efectiva al sembrar
se evalúan caso por caso; el valor derivado \(c\) no aporta entropía adicional.
& Dos bloques, copia al iniciar y XOR con el estado al resembrar. Se exige una cota conjunta.
& Tres palabras, máscara y Update. La proyección a la clave puede costar hasta
128 bits, aun con estado completo de 384.\\
.2, .3 & §§1068--1069. El contador público no añade incertidumbre; BSI relaciona
el estado eficaz con la información que conserva su observación.
& Longitud 128. Separación por firmas de cinco textos bajo un modelo de claves;
dificultad de inferencia sujeta a la primitiva y la semilla.
& Longitud 128. §102 limita el crédito efectivo a la clave. Separación de
clases en modelo ideal; la fuerza de inferencia concreta se analiza aparte.\\
.5, .6 & §§1053--1064. Predecir el sucesor o el predecesor de una ventana conocida.
& Sustitución sucesiva de las parejas AES/SHA de cada petición.
& Sustitución de los tramos de AES; se incluye la actualización entre peticiones.\\
.7, .8 & §§1065--1067. Conocido el estado posterior, BSI deriva (5.79)--(5.80)
y examina el cálculo de la salida pasada. §1070 trata la elección de entrada.
& Se permite que el test vea la clave terminal y que la entrada dependa del historial autorizado.
& Se permite que el test vea la terna terminal desenmascarada. La política
puede omitir o elegir el dato antes de cada petición.\\
.9, .10 & §§1071--1072. Composición criptográfica, apoyada en el análisis
de inversión anterior, y comportamiento de tests justos durante la vida de la semilla.
& Preimagen de SHA para el sufijo dado. El test de salida puede ignorar
el estado terminal. & Terna de cifrados consecutivos con máscara conocida.
El mismo análisis por tramos cubre los tests de salida.\\
\bottomrule
\end{tabularx}
\end{center}

\paragraph{Cómo instanciar la vista de .5, .6 y .7.}
Sea \(Y_1,\ldots,Y_N\) la salida de una época y \(i\le j\).
Para .5 se concede \(W=(Y_i,\ldots,Y_j)\) y se pide \(Y_{j+1}\);
para .6 se concede la misma ventana y se pide \(Y_{i-1}\), cuando existen
los índices. Para .7 se fija el punto de compromiso después de la petición
\(N\) y se incluye en \(W\) el estado vigente y la última entrada si procede.
El árbitro registra los bits objetivo para verificar el acierto al final;
el predictor y la política de entradas sólo leen \(W\) y sus monedas.
Por ejemplo, en .6 el árbitro retiene el bloque anterior y al completar
la ventana evalúa la conjetura calculada a partir de ésta. En el ideal, para
cualquier selección fija de \(t\) bits ocultos, cada valor permitido de \(W\)
con masa positiva admite \(2^t\) valores equiprobables del objetivo:
la ventana usa coordenadas disjuntas, el estado terminal ideal es independiente
y los datos elegidos se computan a partir de la vista autorizada. El acierto
ideal vale \(2^{-t}\). Una elección adaptativa de posiciones requiere verificar
esa igualdad para su política concreta. El índice de la época y las condiciones
de salud forman parte del modelo si se hacen visibles.
Así se aplica §214 también a objetivos que atraviesan una petición.
Las definiciones de observador en Lean admiten árbitros generales; el cálculo
del acierto ideal requiere precisamente esta restricción sobre el predictor.

\paragraph{Entrada controlada e inconspicuidad.}
Para .8, sean \(p_\star(T,t)\) la probabilidad de acierto tolerada para el
tiempo de ataque \(T\) y \(N\) el horizonte elegido. Se debe justificar que
\emph{todas} las políticas permitidas, incluida la entrada vacía, satisfacen
el objetivo en esa vida de la semilla. El análisis de AES/SHA usa el mismo
máximo \(N\) y el de CTR hasta \(2N\) tramos; el modo autónomo CTR usa \(N\).
Los límites uniformes permiten comprobar un objetivo común; no calculan una
igualdad entre las ventajas reales de dos políticas.

Para .10, elíjase un test justo de rechazo \(D(Y^N)\), ejecutable con el
presupuesto previsto y con probabilidad de rechazo a lo sumo \(p_0\) sobre
la secuencia ideal. El árbitro evalúa el test al final e ignora el estado
terminal. Con semilla uniforme, de los teoremas de transferencia de aceptación
se deduce
\[
\begin{aligned}
 \Pr_{\rm AES/SHA}[D=1]&\le\min(1,p_0+N\alpha_{\rm par}(T)),\\
 \Pr_{\rm CTR}[D=1]&\le\min(1,\gamma^J[p_0+J\alpha_4(T)]) .
\end{aligned}
\]
Aquí \(J=N\) en operación autónoma y \(J=2N\) en el perfil mixto,
\(M=2^{128}\) y \(\gamma=M^4/(M)_4\).
La calidad inicial real se introduce con la transferencia de semilla de
cada expediente. Las doce observaciones de la sección~\ref{sec:rates-checked}
comprueban casos reproducibles. La condición sobre cualquier test elegible
proviene de las hipótesis computacionales del expediente. El §200, aplicable
a DRG.3.10 por §209, admite fundamento teórico, fundamento teórico respaldado
por ensayos o ensayos con selección justificada.

\paragraph{Suficiencia de la evidencia.}
Para PTG.3.2, el §380 remite a .3 y .5--.8. Ambos núcleos deben sostener .5--.8
por su operación algorítmica incluso sin contribución de entropía fresca (§211).
En AES/SHA hay un precedente explícito del BSI y una hipótesis conjunta
suficiente para la cota cuantitativa elegida. Para CTR hay pruebas de las
transiciones y reducciones a seguridad PRP, un argumento estructural en el
modelo ideal y una limitación cuantitativa concreta de la reducción por
tramos para .3. En ambos casos quedan por acreditar el soporte de la fuente
real, el horizonte efectivo, los costes de ataque de la instancia y la
correspondencia del TOE con el perfil estudiado. El ejemplo de Hash\_DRBG
identifica de modo análogo las obligaciones de fuente fuera de su conclusión
algorítmica en §§1073--1075.
